% SOURCE_STATEMENT: "Do analogues of Osterwalder–Schrader positivity exist in free probability or for noncommuting random variables?"
% Preparation and final self-audit: CONTRIBUTING.md prompts C and D; source/audit artifacts retained in the local physics_statements_2026_10_11 batch.
% !TeX program = lualatex
\documentclass[11pt]{article}
\usepackage[margin=25mm]{geometry}
\usepackage{fontspec}
\setmainfont{Latin Modern Roman}
\usepackage{amsmath,amssymb,mathtools}
\usepackage{unicode-math}
\setmathfont{Latin Modern Math}
\usepackage{xurl,hyperref}
\hypersetup{colorlinks=true,urlcolor=blue}
\setcounter{secnumdepth}{0}
\setlength{\parindent}{0pt}
\setlength{\parskip}{5pt}
\setlength{\emergencystretch}{3em}
\newcommand{\sref}[2]{\hyperlink{source:#1:#2}{[S#2]}}
\newcommand{\cataloguescope}[1]{\paragraph{Recorded target.}#1}
\newcommand{\R}{\mathbb R}
\newcommand{\C}{\mathbb C}
\newcommand{\N}{\mathbb N}
\newcommand{\Z}{\mathbb Z}
\newcommand{\Q}{\mathbb Q}
\newcommand{\F}{\mathbb F}
\newcommand{\eps}{\varepsilon}
\begin{document}
\section{Osterwalder–Schrader Positivity in Free Probability}\label{statement}

\subsection{Definitions and mathematical statement}

\paragraph{Shared notation.}$\mathbb N=\{1,2,\ldots\}$; the standard number fields have their usual meanings. All additional conventions are stated locally.


\paragraph{Definitions and conventions.}
Free probability studies a noncommutative probability space $(\mathcal A,\varphi)$, with a positive normalized functional $\varphi$ and a notion of free independence. Noncommuting random variables are elements of $\mathcal A$ whose products depend on order. An analogue of reflection positivity requires a reflection or anti-automorphism $\theta$ and a distinguished positive-time algebra, with a positive sesquilinear form corresponding to $\varphi(\theta(a)^*b)$ under the chosen convention.

\paragraph{Target retained from the canonical source catalogue.}
\paragraph{Catalogue wording retained for traceability.}
Do analogues of Osterwalder–Schrader positivity exist in free probability or for noncommuting random variables?

\paragraph{Source scope and qualification.}
The reflection, ordering convention and positive-time algebra are to be constructed or characterized. The displayed form only explains the type of positivity sought; it is not a universal axiom already satisfied by every noncommutative state.
\subsection{Short English statement}
Can Euclidean reflection positivity be formulated meaningfully for free probability and noncommuting random variables?
\subsection{Sources}
\begin{itemize}
\item[\hypertarget{source:statement:1}{S1}] ULAM.AI and the UnsolvedMath Contributors, UnsolvedMath, catalogue entry. Catalogue curation under CC BY 4.0; definitions and English exposition added locally.
\url{https://huggingface.co/datasets/ulamai/UnsolvedMath}.
\item[\hypertarget{source:statement:2}{S2}] Reflection Positivity (2017), Oberwolfach Reports, DOI 10.4171/owr/2017/55, pp. 3263--3343. Primary report accessed 2026-10-11; relevant PDF pages 47, 48.
\url{https://doi.org/10.4171/owr/2017/55}.
\item[\hypertarget{source:statement:3}{S3}] Speicher, Bi-Free Probability and Reflection Positivity (catalogue-listed context; not independently consulted in this preparation).
\url{https://arxiv.org/abs/2312.06813}.
\item[\hypertarget{source:statement:4}{S4}] OWR Report 55/2017 (catalogue-listed context; not independently consulted in this preparation).
\url{https://oa.tib.eu/renate/server/api/core/bitstreams/8df8b5cb-72c1-4ca0-929f-dca03ec9bec0/content}.
\end{itemize}
\label{end:statement}
\end{document}
